% Generated draft: requires mathematical and source review.
\documentclass[11pt]{article}
\usepackage[margin=25mm]{geometry}
\usepackage{fontspec}
\setmainfont{Latin Modern Roman}
\usepackage{amsmath,amssymb,mathtools,unicode-math}
\setmathfont{Latin Modern Math}
\usepackage{enumitem,needspace,xurl,hyperref}
\setlength{\parindent}{0pt}
\setlength{\parskip}{5pt}
\setlength{\emergencystretch}{3em}
\newcommand{\N}{\mathbb{N}}
\newcommand{\Z}{\mathbb{Z}}
\newcommand{\Q}{\mathbb{Q}}
\newcommand{\R}{\mathbb{R}}
\newcommand{\C}{\mathbb{C}}
\newcommand{\F}{\mathbb{F}}
\newcommand{\A}{\mathbb{A}}
\newcommand{\PP}{\mathbb{P}}
\newcommand{\E}{\mathbb{E}}
\newcommand{\eps}{\varepsilon}
\newcommand{\dd}{\,\mathrm{d}}
\begin{document}
\section*{UnsolvedMath ID: 6200076}
\textbf{Source problem code:} catalogue entry
\subsection{Definitions and mathematical statement}
For a metric space $(X,d)$, basepoints $x_j$, scales $r_j\to\infty$, and a nonprincipal ultrafilter $\omega$ on $\N$, an asymptotic cone is the metric ultralimit of $(X,r_j^{-1}d,x_j)$. Concretely, consider sequences $(y_j)$ with $d(y_j,x_j)/r_j$ bounded and identify sequences at zero pseudodistance
\[
d_\omega((y_j),(z_j))=\lim_\omega\frac{d(y_j,z_j)}{r_j}.
\]
The quotient carries the induced metric. Changing basepoints, scales, or ultrafilter may change the cone, so these are part of its data.

The source asks what analytic tools can be developed to study asymptotic cones beyond descriptions up to homeomorphism. This is an open-ended programme, not a proposed inequality or a universal existence theorem. It does not specify a class of original spaces or groups, a canonical measure on their cones, regularity assumptions, or a desired analytic output. For example, differentiation, integration, and function-space methods require additional metric or measure structure and hypotheses that a homeomorphism type does not specify. No particular such structure is asserted to exist on every cone.

The source's statement that earlier study had been non-analytic is retained as historical motivation, not a verified description of all current work.
\subsection{Short English statement}
Develop analysis on asymptotic cones, retaining their metric and any appropriate additional structure rather than only their topology. The source proposes a programme without a single quantitative target.
\subsection{Sources}
\begin{itemize}
\item Local UnsolvedMath registry, permanent ID 6200076, exact source code catalogue entry, ``Boundaries of Groups and Kleinian Groups---Problem 76''; proposed by Mario Bonk, 2005. Registry status open, classification OPEN-TRIAGE.
\item Kapovich, \emph{Problems on Boundaries of Groups and Kleinian Groups} (2005), Problem 76, PDF page 20, \url{https://www.math.ucdavis.edu/~kapovich/EPR/problems.pdf}. Supplied normalized pdf-text extraction accessed July 29, 2026. No PDF was accessed here.
\item The supplied research summary calls this a broad open programme and mentions analytic structure for $\R$-trees/hyperbolic settings without precise tools, hypotheses, or references. This is contextual supplied assessment, not a theorem about all asymptotic cones.
\item Provenance flags current status and statement-reuse terms NEEDS\_REVIEW; the source is publicly author-hosted.
\end{itemize}
\end{document}
